\documentclass[10pt,a4paper]{amsart}
\usepackage[margin=19mm]{geometry}
\usepackage{amsmath,amssymb,mathtools}
\usepackage[scr=rsfs,cal=euler]{mathalfa}
\usepackage{xfrac}
\usepackage[unicode,hidelinks,bookmarks=false]{hyperref}
\newcommand{\ie}{i.e.\ }
\newcommand{\cf}{cf.\ }
\newcommand{\resp}{respectively\ }
\newcommand{\Subsection}{\subsection}
\providecommand{\nobreakdash}{\nobreak\hbox{-}}
\setlength{\emergencystretch}{2em}
\multlinegap0pt
\usepackage{mathtools}
\usepackage{xcolor}
\newtheorem{theorem}{Theorem}
\newcommand{\G}{\xi}
\newcommand{\cfn}{C(\sn )_{\scriptscriptstyle +}}
\renewcommand{\d}{\,\mathrm{d}}
\newcommand{\hm}{\mathcal H^{n-2}}
\newcommand{\hn}{\mathcal H^{n-1}}
\newcommand{\measn}{M(\sn )_{\scriptscriptstyle +}}
\newcommand{\oV}{\operatorname{V}}
\newcommand{\on}{\operatorname{O}(n)}
\newcommand{\pole}{e}
\newcommand{\radon}{\operatorname{\overline R}}
\newcommand{\rhot}{\tau}
\newcommand{\rp}{\mathbb R_{\scriptscriptstyle +}}
\newcommand{\sm}{S^{n-2}}
\newcommand{\sn}{{S^{n-1}}}
\newcommand{\somename}{strong Cara\-th\'eo\-dory and uniformly controlled }
\title{Measure-valued valuations on star bodies}
\author{Jorge S. Ib\'a\~nez-Marcos, Monika Ludwig, Pedro Tradacete, Ignacio Villanueva}
\date{Source: arXiv:2606.30538v1 (CC BY 4.0)}
\begin{document}
\maketitle

\noindent\emph{Source and licence.} Measure-valued valuations on star bodies. Version arXiv:2606.30538v1 (CC BY 4.0), distributed by arXiv under the Creative Commons Attribution 4.0 International licence, http://creativecommons.org/licenses/by/4.0/, as recorded in the arXiv licence metadata for this exact version and shown on the abstract page https://arxiv.org/abs/2606.30538v1. Full licence notice: \texttt{licenses/arxiv-2606.30538-CC-BY-4.0.txt}.\par\smallskip

\noindent\textit{Editorial scope.} Counted unit: the article's theorem equivariant (source0.tex lines 971--977). The article's complete original proof of it is delivered with it (source0.tex lines 978--1042, one \texttt{proof} environment of 65 lines). No mathematics is written, repaired or re-derived: the statement and the proof are the article's own lines, in the article's own notation, and no retained line is substituted or re-flowed.  Retained uncounted as the source-local context the proof needs: lines 680--686; lines 954--956.  Nothing was omitted from any retained span, so the package carries no omission record.  No retained line contains a citation, so the wrapper carries no bibliography and no bibliography entry.  No retained line contains a figure environment, an image-inclusion command or drawing code, so no image is delivered and the package declares no asset dependency.  Editorial formatting only, in addition: an AMS \texttt{amsart} preamble that restates the article's own declarations and macros used by the retained lines, namely the environments \texttt{theorem} on separate counters, together with the standard packages those lines need.  The retained environments print in this document as Theorem 1 in order of appearance, whereas the article prints the same items with the numbering of its own full document.  The complete native source of the article as distributed in the arXiv e-print for this exact version is bundled unchanged as \texttt{sources/204031/source0.tex}.
\begin{theorem}\label{th:representation}
A map $\oV\colon  \cfn\to \measn$ is a weak$^*$\! continuous valuation with $\oV(0)=0$ if and only if there exist $\mu\in \measn$ and a $\mu$-\somename map $U\colon \rp \times \sn  \to \measn$ with $U(0,s)=0$ for $\mu$-almost every $s\in \sn$, such that
\begin{equation}\label{eq_equality}
\langle \oV(f),g\rangle = \int_{\sn } \langle U(f(s), s),g\rangle \d\mu(s)
\end{equation}
for every $f\in\cfn$ and every $g\in C(\sn)$. 
\end{theorem}

\begin{equation}\label{eq:sigma}
\sigma_\lambda \otimes \overline{\mathcal{H}}^{n-2}    + a_\lambda \,\delta_{\pole} + b_\lambda \delta_{-\pole} 
\end{equation}

\begin{theorem}\label{th:representation equivariant}
The map $\oV\colon \cfn\to \measn$ is a weak$^*$~continuous and rotation equivariant valuation with $\oV(0)=0$ if and only if there exists a weak$^*$~continuous map $\rhot\colon \rp\to M([-1,1])_{\scriptscriptstyle +}$  with $\rhot_0=0$ such that
\[
\langle V (f),g\rangle= \int_\sn \int_{[-1,1]}\radon g(s,\alpha)\d\rhot_{f(s)}(\alpha)\d\hn(s) 
\]
for any $f\in \cfn$ and $g\in C(\sn)$.
\end{theorem}

\begin{proof}
First, we will define $\rhot_\lambda$ by duality. For any $\G\in C([-1,1])$ we define
\[
    \rhot_\lambda(\G) =\int_{(-1,1)} \G(\alpha)\d\sigma_\lambda(\alpha)+ a_\lambda     \G(1)+     b_\lambda \G(-1)
\]
using \eqref{eq:sigma}.
Clearly, $\rhot_\lambda$ is linear and bounded and, therefore, $\rhot_\lambda$ is a measure on $[-1,1]$. Moreover,
\begin{align*}
    \langle \rhot_\lambda,\radon g (s, \cdot)\rangle &=\int_{(-1,1)} \left(\int_{\sm}g(\vartheta_s^{-1}(\alpha,t))\d \overline{\mathcal{H}}^{n-2}(t)\right)\d\sigma_\lambda(\alpha)\\
    &\hspace{5mm}+a_\lambda g(s)+b_\lambda g(-s)\\
    &=\langle U'(\lambda,s),g\rangle.
\end{align*}
Now, we have to check that $\rhot_\lambda $ is weak$^*$ continuous but this is a consequence of two facts, the surjectivity of $g\mapsto \radon g (s,\cdot)$ and the weak$^*$ continuity of the mapping $\lambda\mapsto U'(\lambda,s)$. Indeed, let $\G\in C([-1,1])$, $\lambda\in\rp$ and $\{\lambda_m\}_{m\in\mathbb{N}}\subset \rp$ a sequence that converges to $\lambda$. Since $g\mapsto \radon g (s,\cdot)$ is a surjective mapping, there exists $g\in C(\sn)$ such that $\radon g(s, \cdot)=\G$ and, thus,
\[
    \langle \rhot_{\lambda_m},\G\rangle=\langle \rhot_{\lambda_m},\radon g (s, \cdot)\rangle =\langle U'(\lambda_m,s),g\rangle\to \langle U'(\lambda,s),g\rangle=\langle \rhot_\lambda,\G\rangle
\]
as $m\to\infty$.

\goodbreak
Conversely, let $\rhot_\lambda\in M([-1,1])_{\scriptscriptstyle +}$. For any $\lambda\in\rp$ and any $s\in\sn$, define the mapping
\[
     U'(\lambda,s)(g) =\int_{[-1,1]} \radon g (s, \alpha)\d \rhot_\lambda (\alpha)=\langle \rhot_\lambda, \radon g (s, \cdot)\rangle.
\]
Clearly, this mapping is linear and continuous; therefore, it defines a measure. Observe that that $\|U'(\lambda,s)\|\leq \|\rhot_\lambda\|$; hence, it is uniformly controlled. We now verify that it is $\mathcal{H}^{n-1}$-strong Carath\'eodory. First, it follows easily that the mapping
\[
    \lambda\mapsto \langle U'(\lambda,s),g\rangle
\]
is continuous for any $g\in C(\sn)$ and any $s\in \sn$. For any $\alpha\in(-1,1)$ and $e\in\sn$, define
\[
    C_\alpha=
    \frac{1}{\hm(\{s\in\sn:\langle s,e\rangle=\alpha\})}.
\] 
Fix $g\in C(\sn)$ and $\lambda\in\rp$. Let $s\in \sn$ and let $\{s_m\}_{m\in\mathbb{N}}\subset\sn$ be a sequence that converges to $s$. Then
\[
    \radon g (s_m, \pm 1)=g(\pm s_m)
\]
and, for $\alpha\in(-1,1)$,
\begin{align*}
   \radon g (s_m, \alpha)&=\int_{\sm}g(\vartheta^{-1}_{s_m}(\alpha,t))\d\overline{\mathcal{H}}^{n-2}(t)\\
    &=C_\alpha \int_{\{r\in\sn:\langle r,s_m\rangle=\alpha\}}g(r)\d\hm(r).
\end{align*} 
Since $g$ is uniformly continuous, it follows that $\radon g (s_m, \cdot)$ converges uniformly to $\radon g (s, \cdot)$. Hence, the mapping 
\[
    s\mapsto \langle U'(\lambda,s),g\rangle
\]
is continuous and, in particular, it is Borel measurable. Thus, by Theorem \ref{th:representation}, the mapping $f\mapsto \oV (f)$, given by
\begin{align*}
    \langle \oV(f),g\rangle &=\int_\sn \langle U'(\lambda,s),g\rangle \d \hn(s)\\
    &=\int_{\sn}\int_{[-1,1]} \radon g (s, \alpha)\d\rhot_\lambda(\alpha)\d\hn(s),
\end{align*}
is a weak$^*$ continuous valuation.

Finally, let us check that $\oV$ is rotation equivariant. Note that it is enough to see that $\radon (\phi g) (\phi s, \cdot)= \radon g (s, \cdot)$ for any $g\in C(\sn)$ and any $\phi\in\on$. This follows easily since
\[
   \radon (\phi g) (\phi s, \pm 1) =\phi g(\phi(\pm s))=g(\pm s),
\]
and for any $\alpha\in(-1,1)$ we have
\begin{align*}
   \radon (\phi g) (\phi s, \alpha)&=\int_{\sn}\phi g(\vartheta_{\phi(s)}^{-1}(\alpha,t))\d \overline{\mathcal{H}}^{n-2}(t)\\
    &=C_\alpha \int_{\{r\in\sn: \langle r,\phi(s)\rangle=\alpha\}}\phi g(r)\d \hm(r)\\
    % &=C_\alpha \int_{\{r\in\sn:\langle \phi(r),\phi(s)\rangle =\alpha\} }g(r)\d\hm(r)\\
    &=C_\alpha \int_{\{r\in\sn:\langle s,r\rangle =\alpha\} }g(r)\d\hm(r)=\radon g (s, \alpha).
\end{align*}
Therefore, the valuation $\oV$ is rotation equivariant.
\end{proof}

\end{document}
